\documentclass[12pt,letterpaper]{article}
\usepackage{graphicx,textcomp}
\usepackage{natbib}
\usepackage{setspace}
\usepackage{fullpage}
\usepackage{color}
\usepackage[reqno]{amsmath}
\usepackage{amsthm}
\usepackage{fancyvrb}
\usepackage{amssymb,enumerate}
\usepackage[all]{xy}
\usepackage{endnotes}
\usepackage{lscape}
\newtheorem{com}{Comment}
\usepackage{float}
\usepackage{hyperref}
\newtheorem{lem} {Lemma}
\newtheorem{prop}{Proposition}
\newtheorem{thm}{Theorem}
\newtheorem{defn}{Definition}
\newtheorem{cor}{Corollary}
\newtheorem{obs}{Observation}
\usepackage[compact]{titlesec}
\usepackage{dcolumn}
\usepackage{tikz}
\usetikzlibrary{arrows}
\usepackage{multirow}
\usepackage{xcolor}
\newcolumntype{.}{D{.}{.}{-1}}
\newcolumntype{d}[1]{D{.}{.}{#1}}
\definecolor{light-gray}{gray}{0.65}
\usepackage{url}
\usepackage{listings}
\usepackage{color}

\definecolor{codegreen}{rgb}{0,0.6,0}
\definecolor{codegray}{rgb}{0.5,0.5,0.5}
\definecolor{codepurple}{rgb}{0.58,0,0.82}
\definecolor{backcolour}{rgb}{0.95,0.95,0.92}

\lstdefinestyle{mystyle}{
	backgroundcolor=\color{backcolour},   
	commentstyle=\color{codegreen},
	keywordstyle=\color{magenta},
	numberstyle=\tiny\color{codegray},
	stringstyle=\color{codepurple},
	basicstyle=\footnotesize,
	breakatwhitespace=false,         
	breaklines=true,                 
	captionpos=b,                    
	keepspaces=true,                 
	numbers=left,                    
	numbersep=5pt,                  
	showspaces=false,                
	showstringspaces=false,
	showtabs=false,                  
	tabsize=2
}
\lstset{style=mystyle}
\newcommand{\Sref}[1]{Section~\ref{#1}}
\newtheorem{hyp}{Hypothesis}

\title{Problem Set 2}
\date{Due: October 22, 2026}
\author{Ethan Kelly Zeh (26359218)}

\begin{document}
	\maketitle
	\section*{Instructions}
\begin{itemize}
	\item Please show your work! You may lose points by simply writing in the answer. If the problem requires you to execute commands in \texttt{R}, please include the code you used to get your answers. Please also include the \texttt{.R} file that contains your code. If you are not sure if work needs to be shown for a particular problem, please ask.
	\item Your homework should be submitted electronically on GitHub.
	\item This problem set is due before 23:59 on Thursday October 22, 2026. No late assignments will be accepted.

\end{itemize}

	
	\vspace{.5cm}
	\section*{Question 1: Political Science}
		\vspace{.25cm}
	The following table was created using the data from a study run in a major Latin American city.\footnote{Fried, Lagunes, and Venkataramani (2010). ``Corruption and Inequality at the Crossroad: A Multimethod Study of Bribery and Discrimination in Latin America. \textit{Latin American Research Review}. 45 (1): 76-97.} As part of the experimental treatment in the study, one employee of the research team was chosen to make illegal left turns across traffic to draw the attention of the police officers on shift. Two employee drivers were upper class, two were lower class drivers, and the identity of the driver was randomly assigned per encounter. The researchers were interested in whether officers were more or less likely to solicit a bribe from drivers depending on their class (officers use phrases like, ``We can solve this the easy way'' to draw a bribe). The table below shows the resulting data.

\newpage
\begin{table}[h!]
	\centering
	\begin{tabular}{l | c c c }
		& Not Stopped & Bribe requested & Stopped/given warning \\
		\\[-1.8ex] 
		\hline \\[-1.8ex]
		Upper class & 14 & 6 & 7 \\
		Lower class & 7 & 7 & 1 \\
		\hline
	\end{tabular}
\end{table}

\begin{enumerate}
	
	\item [(a)]
	Calculate the $\chi^2$ test statistic by hand/manually (even better if you can do "by hand" in \texttt{R}).\\
	\vspace{2cm}

	\textbf{Explanation}
	\vspace{.5cm}

	The following code manually calculates the Chi-Squared ($\chi^2$) test statistic using the formula:
	\[
	\chi^2 = \sum_{i=1}^{n} \frac{(F_o - F_e)^2}{F_e}
	\]
	The code first creates the $2 \times 3$ matrix. Following this, we have two nested for-loops, which iterate by the column number and row number, therefore allowing us to do the calculation at each index in the matrix. Lines 12 and 13 calculate the total at that row and column respectively. Line 14 calculates the expected value $F_e$. Line 16 calculates the summation at $n$, which, on line 17, we add to the running total variable \texttt{total\_sum\_psm}. Ultimately, we assign this to the test statistic variable ($\chi^2$) \texttt{poly\_sci\_test\_stat}. The value is printed and the output is: 3.791168

	This manual approach was done in a programmatic way. This means we could change the matrix/values and it would still work as intended, despite not using R's internal Chi-Squared functions.
	\vspace{1cm}

	\textbf{Code}
    \lstinputlisting[language=R, firstline=2, lastline=22]{PS02_answersEZ.R}
	\vspace{1cm}
    \textbf{Output}
	\begin{verbatim}
        3.791168
	\end{verbatim}

	
	\item [(b)]
	Now calculate the p-value from the test statistic you just created (in \texttt{R}).\footnote{Remember frequency should be $>$ 5 for all cells, but let's calculate the p-value here anyway.}  What do you conclude if $\alpha = 0.1$?\\
	
	\vspace{2cm}

	\textbf{Explanation}
	\vspace{.5cm}

	The following are our null and alternate hypothesis:
	\begin{align*}
	H_0 &: \text{Driver class and officer response are independent.} \\
	H_a &: \text{Driver class and officer response are not independent.}
	\end{align*}

	Using a right-tailed Pchisq method with 3 degrees of freedom (as seen in code below), we calculate the P value at approx. 0.15. Because $p = 0.15 > 0.10$, we fail to reject the null hypothesis. There is insufficient evidence at the $\alpha = 0.10$ level to conclude that the likelihood of officers soliciting a bribe differs by driver class. This does not prove that these variables are indepedent, but rather that this study does not show sufficient proof otherwise. The small sample size is a major factor driving this outcome. As mentioned in the footnote, there should be at least 5 samples in each cell.
	\vspace{1cm}

	\textbf{Code}
    \lstinputlisting[language=R, firstline=25, lastline=27]{PS02_answersEZ.R}
	\vspace{1cm}
    \textbf{Output}
	\begin{verbatim}
        0.1502306
	\end{verbatim}
	
	\newpage
	\item [(c)] Calculate the standardized residuals for each cell and put them in the table below.
	\vspace{1cm}
	
	\begin{table}[h]
		\centering
		\begin{tabular}{l | c c c }
			& Not Stopped & Bribe requested & Stopped/given warning \\
			\\[-1.8ex] 
			\hline \\[-1.8ex]
			Upper class  &0.322  &-1.642  &1.523 \\
			\\
			Lower class &-0.322  &1.642   &-1.523  \\
			
		\end{tabular}
	\end{table}
	\vspace{1cm}
	The results for this table were calculated using the following formula manually in R, similar to the programatic approach of question 1. a
	\vspace{.5cm}
		\[
		z_{ij} = \frac{F_o - F_e}{\sqrt{F_e\left(1 - \dfrac{n_{i\cdot}}{n_{ij}}\right)\left(1 - \dfrac{n_{\cdot j}}{n_{ij}}\right)}}
		\]
		Where:
		\vspace{.5cm}
		\[ 
		F_e = \frac{n_{i\cdot}\, n_{\cdot j}}{n_{ij}}
		\] 
	\vspace{.5cm}

	The code iterates through each column and then row, grabs the total for that row and column, calculates $F_e$ as well as the standard error (the denominator in the above equation), calculates the std. residual, and then prints it in Table form. The output of the code is below.
	\vspace{.5cm}


	\textbf{Code}
    \lstinputlisting[language=R, firstline=29, lastline=48]{PS02_answersEZ.R}
	\vspace{1cm}
    \textbf{Output}
	\begin{verbatim}

	 Row: 1 | Column: 1 | Value: 14.0000 | Z (std res.): 0.3220
		Row: 2 | Column: 1 | Value: 7.0000 | Z (std res.): -0.3220
		Row: 1 | Column: 2 | Value: 6.0000 | Z (std res.): -1.6420
		Row: 2 | Column: 2 | Value: 7.0000 | Z (std res.): 1.6420
		Row: 1 | Column: 3 | Value: 7.0000 | Z (std res.): 1.5230
		Row: 2 | Column: 3 | Value: 1.0000 | Z (std res.): -1.5230
	\end{verbatim}
	
	
	\vspace{2cm}
	\item [(d)] How might the standardized residuals help you interpret the results?  
	
	\vspace{1cm}
	Lower-class drivers were asked for bribes more often than expected under independence (residual $+1.642$) and upper-class drivers less often ($-1.642$), but no residual is greater than 2. Therefore, no cell departs significantly from independence. This aligns with our p-value of .15, as we do not see significant results.
	
\end{enumerate}
\newpage

\section*{Question 2: Economics}
Chattopadhyay and Duflo were interested in whether women promote different policies than men.\footnote{Chattopadhyay and Duflo. (2004). ``Women as Policy Makers: Evidence from a Randomized Policy Experiment in India. \textit{Econometrica}. 72 (5), 1409-1443.} Answering this question with observational data is pretty difficult due to potential confounding problems (e.g. the districts that choose female politicians are likely to systematically differ in other aspects too). Hence, they exploit a randomized policy experiment in India, where since the mid-1990s, $\frac{1}{3}$ of village council heads have been randomly reserved for women. A subset of the data from West Bengal can be found at the following link: \url{https://raw.githubusercontent.com/kosukeimai/qss/master/PREDICTION/women.csv}\\

\noindent Each observation in the data set represents a village and there are two villages associated with one GP (i.e. a level of government is called "GP"). The authors hypothesize that female politicians are more likely to support policies female voters want. Researchers found that more women complain about the quality of drinking water than men. You need to estimate the effect of the reservation policy on the number of new or repaired drinking water facilities in the villages.
\vspace{.5cm}


\newpage
\begin{enumerate}
	\item [(a)] State a null and alternative (two-tailed) hypothesis. 
	\vspace{.5cm}

	The following are our null and alternate hypothesis:
	\begin{align*}
	H_0 &: \text{Reservation policy has no effect on the number of new or repaired drinking water facilities.} \\
	H_a &: \text{Reservation policy changes the number of new or repaired drinking water facilities.}
	\end{align*}

	\vspace{2cm}
	\item [(b)] Run a bivariate regression to test this hypothesis in \texttt{R} (include your code!).
	\vspace{.5cm}
		
	\textbf{Explanation}
	\vspace{.5cm}

	I used the following formula to manually calculate the coeffecient estimate (slope) of this relationship. I chose the \texttt{water} and \texttt{reserved} variables because I wanted to see how the policy of gender-quota affects the outcome, not solely women as council-heads.

	\[
	{\beta} = \frac{\sum_{i=1}^{n}(x_i - \bar{x})(y_i - \bar{y})}{\sum_{i=1}^{n}(x_i - \bar{x})^2}
	\]
	\vspace{.5cm}
	
	I then compared this outcome (approx. 9.25) with the outcome of the \texttt{lm} function in R in order to confirm the coeffecient estimate and find the- p-value.
	
	\vspace{.5cm}

	\textbf{Code}
    \lstinputlisting[language=R, firstline=55, lastline=71]{PS02_answersEZ.R}
	\vspace{1cm}
    \textbf{Output}
	\begin{verbatim}
        9.252423

		Coefficients:
            Estimate Std. Error t value Pr(>|t|)    
(Intercept)   14.738      2.286   6.446 4.22e-10 ***
reserved       9.252      3.948   2.344   0.0197 * 
	\end{verbatim}
	\vspace{2cm}
	\item [(c)] Interpret the coefficient estimate for reservation policy. 
	\vspace{1cm}

The estimated coefficient on \texttt{reserved} is $9.252$, meaning villages with a council head reserved for a woman had on average about 9.25 more new or repaired drinking water facilties than unreserved villages. Since $p = 0.0197 < 0.05$, we reject the null hypothesis, and because reservation was randomly assigned, this difference could be interpreted as a causal effect of the policy. Of note: The residuals are strongly right-skewed (minimum $-23.991$, maximum $316.009$), possibly suggesting that a few villages with very large numbers of new or repaired facilities may be driving the estimate. Further examination of this to check for robustness to these outliers could be valuable (ie. dropping extreme values).\end{enumerate}


%\end{enumerate}
\end{document}
